\graphicspath{{abstracts/praesentationen/03-iiif-vernetzung/img/}{img/}}

\begin{beitrag}{IIIF und WissKI: Vernetzung digitaler Bildrepositorien für die Sammlungsforschung}{%
Elena Engel\,\orcidlink{0000-0002-6789-0123} \\
\small \href{mailto:elena.engel@example.org}{elena.engel@example.org} \\
\small Universität Wien \\[0.5cm]
Frank Förster\,\orcidlink{0000-0003-7890-1234} \\
\small \href{mailto:frank.foerster@example.org}{frank.foerster@example.org} \\
\small Universität zu Köln
}

\section{Einleitung}
Das \emph{International Image Interoperability Framework} (IIIF) hat sich in den vergangenen Jahren als de-facto-Standard für den interoperablen Austausch digitaler Bildressourcen etabliert. Sammlungseinrichtungen stellen ihre Bestände zunehmend über IIIF-konforme Bildserver bereit, während die Forschung darauf aufbauend annotierende, vergleichende und kuratierende Werkzeuge entwickelt. Vor diesem Hintergrund stellt sich für WissKI-basierte Sammlungsprojekte die Frage, wie IIIF-Ressourcen aus externen Repositorien systematisch in das eigene semantische Datenmodell eingebunden werden können -- nicht als bloße Verlinkung, sondern als integraler Bestandteil der Forschungsdaten. Der Beitrag stellt einen prototypischen Workflow vor, der im Kontext eines kunsthistorischen Forschungsprojekts entwickelt und über zwei Jahre produktiv erprobt wurde.

\section{Ansatz}
Im Zentrum des Workflows steht eine schmale Vermittlungsschicht zwischen IIIF Presentation API und WissKI. Eingangsseitig werden IIIF-Manifeste analysiert, ihre Strukturmetadaten (Manifest, Range, Canvas, Annotation) extrahiert und auf eine an CIDOC CRM angelehnte Klassenhierarchie abgebildet. Ausgangsseitig erzeugt die Schicht aus WissKI-Entitäten dynamisch IIIF-konforme Manifeste, sodass die in WissKI erfassten Forschungsdaten auch außerhalb der eigenen Infrastruktur konsumierbar bleiben.

Die Eingangsanalyse hat sich als methodisch deutlich anspruchsvoller erwiesen als zunächst angenommen. IIIF-Manifeste werden in der Praxis sehr heterogen gestaltet: Beschriftungen liegen in unterschiedlichen Sprachen vor, Provenienzfelder werden uneinheitlich befüllt, und Strukturhierarchien (\texttt{Range}) werden häufig nur sehr partiell gepflegt. Wir haben deshalb für jede angebundene Quelle ein leichtgewichtiges Profil definiert, das jene Felder beschreibt, die im Zielsystem benötigt werden, und das die Toleranz gegenüber fehlenden Werten explizit dokumentiert. Eine ausführliche Diskussion der Heterogenitätsproblematik findet sich bei \cite{engel2024-iiif-heterogen}.

\section{Beispielanwendung}
Die im Projekt umgesetzte Beispielanwendung verknüpft Werke aus drei verschiedenen Quellen: einer regionalen Galerie-Datenbank, einem universitären Bildarchiv und einer großen internationalen Sammlung. Forschende können in der WissKI-Oberfläche Bildwerke recherchieren, ihre Annotationen anreichern und über ein Lightbox-ähnliches Werkzeug Bildvergleiche durchführen, ohne dass die zugrunde liegenden Bilddateien lokal vorgehalten werden müssen. Über die IIIF Image API werden Ausschnitte, Skalierungen und Drehungen on demand vom jeweiligen Bildserver geliefert.

Eine zentrale Designentscheidung war es, in WissKI keine eigene Bildhaltung aufzubauen, sondern strikt auf die Adressierung externer IIIF-Ressourcen über URIs zu setzen. Dies hat zwei Konsequenzen: Erstens wird die Verfügbarkeit der Forschungsumgebung an die Verfügbarkeit der externen Bildserver gekoppelt; zweitens entsteht eine vergleichsweise schmale, aber semantisch reichhaltige Datenbasis, die zuverlässig zitierbar und reproduzierbar ist.

\section{Diskussion}
Im Betrieb haben sich zwei Themen als besonders sensibel erwiesen. Erstens die Langzeitstabilität externer IIIF-Endpunkte: Mehrfach haben Quellen ihre Manifest-URLs während des Projektzeitraums geändert, was zu temporären Ausfällen geführt hat. Wir empfehlen daher den Aufbau einer projektseitigen Resolver-Schicht, die bekannte Verschiebungen abbildet und im Notfall auf gespiegelte Manifest-Kopien zurückgreift. Zweitens die Frage der Lizenzierung: Auch wenn IIIF-Manifeste grundsätzlich offen ausgeliefert werden, sind die zugehörigen Bilddateien häufig mit Nutzungseinschränkungen versehen, die in der WissKI-Anwendung adäquat dargestellt werden müssen; vgl.\ \cite{foerster2025-iiif-lizenz}.

\section{Fazit und Ausblick}
Die Integration von IIIF in WissKI ist auch mit überschaubarem Aufwand möglich und eröffnet erhebliche methodische Mehrwerte für sammlungsbasierte Forschungsprojekte. Die im Vortrag vorgestellte Architektur wird im Rahmen einer Open-Source-Veröffentlichung im laufenden Jahr der Community zugänglich gemacht.

\end{beitrag}
